\documentclass{situationist-preprint}

\title{What Constructivism Constructs}
\msid{MTE-18-a}
\msfamily{MTE-18}
\msdiscipline{Metaethics}
\mstype{Principal manuscript}
\mskeywords{constructivism; realism; procedures; original position; pure procedural justice; constitutivism; idealized endorsement}

\begin{document}
\maketitle

\begin{abstract}
Constructivists hold that moral truths are the outputs of a procedure: the principles that would be chosen in the original position, the maxims that pass the categorical imperative test, the principles no one could reasonably reject. Critics reply that constructivism must either say what makes the procedure correct, in which case it is a form of realism, or not, in which case it is arbitrary. We argue that constructivism is a two-level view, and that it is stable. It constructs outcomes, about which it is anti-realist: no output of the correct procedure is true independently of the procedure. It does not construct the procedure, about which it is realist: the correctness of the procedure is not the output of that procedure or of any other. And the procedure is mandatory: every rational agent is required to accept its outputs, whether or not she accepts the procedure. We derive these claims from four premises about construction, authority, transmission and the symmetry of procedural validation, without assuming that grounding is well-founded. It follows that global constructivism is incoherent, and that a constructivism whose procedure is relative to each agent's standpoint collapses into relativism. We illustrate the structure with the original position of POL-01 and the idealized reflective endorsement of ETH-01.
\end{abstract}

\section{The question}

Constructivism entered contemporary metaethics with \citet{Rawls1980}, who presented his theory of justice as an instance of \enquote{Kantian constructivism}: the principles of justice are not truths about an independent moral order, which reason discovers, but the principles that would be agreed to by persons who are free and equal under conditions that represent them as such. The procedure of the original position \citep{Rawls1971} is laid out, or constructed, and the principles are its outputs. \citet{ONeill1989} read Kant's own ethics in the same way, as the construction of principles that could be adopted by a plurality of agents without presupposing any authority beyond reason. \citet{Korsgaard1996} grounded the construction in the standpoint of practical reason: the normativity of moral requirements derives from the reflective structure of agency and from the value an agent must place on her own humanity. \citet{Scanlon1998} gave the contractualist version: an act is wrong if it would be disallowed by any set of principles for the general regulation of behaviour that no one could reasonably reject. \citet{Habermas1990} proposed a discourse ethics on which valid norms are those that could meet with the agreement of all affected as participants in practical discourse.

These views differ in their procedures but agree on a structural claim. There is a procedure, and the moral truths, or some important class of them, are what the procedure yields; they are not truths that the procedure merely detects. \citet[\S14]{Rawls1971} drew the distinction that marks the view. In \emph{perfect} procedural justice there is an independent criterion for the right outcome and a procedure guaranteed to reach it, as when the person who cuts a cake takes the last piece. In \emph{pure} procedural justice there is no independent criterion; the outcome is fair whatever it is, provided the procedure is fair and has been properly followed, as in a fair lottery. Constructivism is the view that morality, or justice, is a matter of pure procedural justice at the level of principles.

The question we address is what, on such a view, is constructed. The critics' challenge is a dilemma. If the constructivist says what makes her procedure the correct one, she appeals to a standard the procedure answers to, and her view is a form of realism, with the procedure as an epistemic device \citep{ShaferLandau2003,HussainShah2006}. If she does not, the choice of procedure is arbitrary, and so are its outputs. \citet{Enoch2009} pressed the dilemma against constructivism about practical reason in general: a constructivism that is global, about all normative truths, cannot say what makes its procedure authoritative without either presupposing an unconstructed normative truth or regressing. \citet{Street2008,Street2010} accepted the global aspiration and drew a relativist conclusion: the procedure is fixed by the practical standpoint of each valuing agent, and what follows from it varies with the agent's values.

We argue that the dilemma's first horn is not a refutation but a description. Constructivism must be realist about its procedure, and it can be so without being realist about its outcomes; and the realism it requires is specifically the claim that the procedure is mandatory for every rational agent. Street's alternative is the view that gives up the mandatoriness of the procedure, and it is relativism.

\section{Framework}\label{s:framework}

\begin{definition}[Procedures and outputs]\label{d:procedure}
A \emph{procedure} is a specified method of arriving at principles or verdicts: a choice situation together with a decision rule, a test applied to maxims, a standard of reasonable rejection, or an idealization of reactions. A claim is an \emph{output} of a procedure if the procedure, correctly applied, yields it.
\end{definition}

\begin{definition}[Constructivism]\label{d:constructivism}
\emph{Constructivism} about a normative domain is the view that there is a procedure $P^*$ such that the truths of the domain are the outputs of $P^*$, and each such truth holds in virtue of being an output of $P^*$.
\end{definition}

\begin{definition}[Realism about a fact]\label{d:realism}
A view is \emph{realist} about a normative fact if it holds that the fact obtains and does not hold in virtue of being an output of any procedure. It is \emph{anti-realist} about a class of normative facts if it holds that each of them holds in virtue of being an output of some procedure.
\end{definition}

Definition~\ref{d:realism} uses \enquote{realism} in the sense at issue between constructivists and their critics: independence from construction. It does not deny the existence of the facts constructed. On a constructivism of the kind defended here there is a fact about which principles the parties in the original position would choose, and it is a natural fact; MTE-01, on which moral facts are natural facts about idealized reactions, is compatible with anti-realism about outcomes in the present sense. What it is incompatible with is the denial that there are such facts.

\begin{definition}[Mandatory procedure]\label{d:mandatory}
A procedure is \emph{mandatory} for an agent if the agent is required to accept its outputs whether or not she accepts the procedure.
\end{definition}

The argument uses five assumptions.

\begin{assumption}[Construction]\label{a:construction}
Constructivism holds of morality, with procedure $P^*$, and $P^*$ has at least one output.
\end{assumption}

\begin{assumption}[Authority]\label{a:authority}
Every moral truth binds every rational agent: each rational agent is required to accept it.
\end{assumption}

\begin{assumption}[Transmission]\label{a:transmission}
An output of a procedure binds an agent only in virtue of the agent's being required to accept that procedure's outputs.
\end{assumption}

\begin{assumption}[Consistency of requirement]\label{a:consistency}
No agent is required to accept each of two incompatible claims.
\end{assumption}

\begin{assumption}[Symmetry of validation]\label{a:symmetry}
For every procedure $P$ there is a rival procedure $P'$, some of whose outputs are incompatible with outputs of $P$, such that for every chain of procedures each of which outputs that the next is to be accepted and which ends in $P$, there is a chain of the same form ending in $P'$.
\end{assumption}

We also use the standard assumption that grounding is irreflexive and transitive, and so asymmetric \citep{Fine2012,Rosen2010}. We do not assume that it is well-founded; MET-10 argues that there is no fundamental level, and the argument below does not depend on denying it.

Assumption~\ref{a:construction} is the constructivist's thesis, taken for the sake of the argument. Assumption~\ref{a:authority} is the claim, which constructivists share with their critics, that morality is not optional; MTE-02 argues that moral truths hold for all rational agents. Assumption~\ref{a:transmission} states what it is for a truth to be constructed rather than detected. If an output bound an agent independently of her being required to accept the procedure's outputs, its authority would come from somewhere other than the procedure, and the procedure would be idle in explaining it; that is the realist's view of the procedure as a detection device. Assumption~\ref{a:consistency} is a minimal coherence condition, and is entailed by ETH-15's thesis that there are no genuine moral dilemmas. Assumption~\ref{a:symmetry} records a familiar fact about procedural validation: a procedure can be written to endorse any procedure, itself included, and for any chain of endorsements there is a parallel chain endorsing a rival. A procedure that outputs that it is itself to be accepted is matched by a rival procedure that outputs the same about itself.

\section{The argument}\label{s:argument}

\begin{lemma}[Anti-realism about outcomes]\label{l:outcome}
Constructivism is anti-realist about every moral truth that is an output of $P^*$.
\end{lemma}

\begin{proof}
By Definition~\ref{d:constructivism}, each such truth holds in virtue of being an output of $P^*$; by Definition~\ref{d:realism}, that is anti-realism about it.
\end{proof}

\begin{proposition}[Mandatoriness]\label{p:mandatory}
$P^*$ is mandatory for every rational agent.
\end{proposition}

\begin{proof}
Let $x$ be a rational agent and $\varphi$ an output of $P^*$, which exists by Assumption~\ref{a:construction}. By Definition~\ref{d:constructivism} and Assumption~\ref{a:authority}, $\varphi$ binds $x$. By Assumption~\ref{a:transmission}, it does so in virtue of $x$'s being required to accept $P^*$'s outputs, and so $x$ is so required. Since Assumption~\ref{a:authority} applies to every rational agent, whatever procedure she accepts, $x$ is required to accept $P^*$'s outputs whether or not she accepts $P^*$.
\end{proof}

\begin{lemma}[No self-construction]\label{l:self}
The fact that $P^*$ is mandatory does not hold in virtue of being an output of $P^*$.
\end{lemma}

\begin{proof}
Write $M$ for the fact that $P^*$ is mandatory. Suppose $M$ held in virtue of being an output of $P^*$. An output holds as a moral truth only if it binds (Assumption~\ref{a:authority}), and it binds only in virtue of $M$ (Assumption~\ref{a:transmission}). So $M$ would hold in virtue of a fact, its being a binding output of $P^*$, that holds in virtue of $M$; by transitivity $M$ would hold in virtue of itself, contrary to irreflexivity.
\end{proof}

\begin{lemma}[No construction all the way]\label{l:external}
Either the fact that $P^*$ is mandatory is not constructed, or it holds in virtue of the output of a finite chain of procedures whose first member's mandatoriness is not constructed.
\end{lemma}

\begin{proof}
By Lemma~\ref{l:self}, if $M$ is constructed, the procedure whose output it is is not $P^*$. Suppose $M$ held in virtue of being an output of a procedure $Q_1$. By Assumption~\ref{a:transmission}, the output binds only in virtue of the fact that $Q_1$ is mandatory, and if that fact is in turn constructed, it holds in virtue of being an output of a procedure $Q_2$, and so on. The chain is either finite, ending in a procedure whose mandatoriness is not constructed, which is the second alternative of the lemma, or infinite, with every link constructed. Since grounding need not be well-founded, the infinite case is not excluded by its length alone, and we exclude it otherwise. By Assumption~\ref{a:symmetry}, there is a rival $P'$ of $P^*$ with incompatible outputs and an infinite chain of the same form ending in $P'$. If chains of this form, with every link constructed, sufficed to make the procedure at their end mandatory, then $P'$ would be mandatory as well, and every rational agent would be required to accept incompatible outputs of $P^*$ and $P'$, contrary to Assumption~\ref{a:consistency}. So they do not suffice, and the infinite case does not arise.
\end{proof}

\begin{theorem}[What constructivism constructs]\label{t:main}
Given Assumptions~\ref{a:construction} to~\ref{a:symmetry}, constructivism is anti-realist about the outputs of its procedure and realist about procedures: some fact of the form \enquote{procedure $Q$ is mandatory} obtains without holding in virtue of being an output of any procedure, and the mandatoriness of $P^*$ either is such a fact or holds in virtue of one. The procedure $P^*$ is mandatory for every rational agent.
\end{theorem}

\begin{proof}
Lemma~\ref{l:outcome}, Proposition~\ref{p:mandatory} and Lemma~\ref{l:external}, with Definition~\ref{d:realism}.
\end{proof}

The theorem answers the question in the title. Constructivism constructs outcomes. It also, in Rawls's phrase, \enquote{lays out} a procedure, and in that sense it presents a construction of procedures as well; but the correctness of the procedure so laid out is not constructed, and constructivism is realist about it. Nothing is gained by placing the unconstructed fact further back, at the head of a chain of procedures each endorsing the next: the procedure at the head is then the one about which the view is realist, and the simplest constructivism, which we assume in what follows, takes the unconstructed fact to be the mandatoriness of $P^*$ itself.

\begin{corollary}[Against global constructivism]\label{c:global}
Not every normative fact is constructed.
\end{corollary}

\begin{proof}
By Theorem~\ref{t:main}, some fact that a procedure is mandatory is not constructed, and such a fact is normative, since it is the fact that every rational agent is required to accept certain outputs.
\end{proof}

\begin{corollary}[Standpoint-relative procedures]\label{c:street}
A constructivism on which the procedure that binds each agent is fixed by that agent's own evaluative standpoint, and on which agents' standpoints may differ, must deny Assumption~\ref{a:authority}, and is a form of relativism.
\end{corollary}

\begin{proof}
If the procedure binding $x$ is fixed by $x$'s standpoint, it is not mandatory for an agent whose standpoint fixes another procedure with incompatible outputs, and by Proposition~\ref{p:mandatory} and Assumption~\ref{a:consistency} that is inconsistent with Assumption~\ref{a:authority}. The truth of a moral claim for an agent then depends on the agent's standpoint, which is relativism in the sense refuted in MTE-06.
\end{proof}

Corollary~\ref{c:street} applies to the Humean constructivism of \citet{Street2008,Street2010}, who accepts its conclusion for the case of an ideally coherent agent whose values include torturing others for fun: such an agent, on her view, may have no reason to refrain. We take this to show that the mandatoriness of the procedure is not an optional addition to constructivism but the condition of its giving morality any authority over those who reject it.

\section{Two constructions}\label{s:cases}

\paragraph{The original position.} POL-01-a defends justice as fairness: the principles of justice for the basic structure are those that would be chosen behind a veil of ignorance, and they include a modified difference principle on which ties on the worst-off level are broken by the total. The principles are outputs of the procedure, and POL-01-a is anti-realist about them in the present sense: there is no fact about the justice of the basic structure that the original position detects. The procedure itself is argued for on other grounds. POL-01-a defends the thick veil of ignorance by appeal to the aim of the device, which is to identify principles that can be justified to every reasonable person subject to them, as POL-02 requires of a legitimate state; and it replies to the objection that the procedure contains in effect a single chooser by saying that the agreement is unanimity, following from the fact that the same reasoning would be accepted by everyone who reasoned under fair conditions. None of these claims is an output of the original position. They are claims about which procedure is correct, made on the basis of a conception of persons as free and equal, and they are the realist component of the view. POL-01-a also illustrates mandatoriness. A person who rejects the veil of ignorance is not thereby released from the principles of justice; the principles bind her because the procedure is the correct one, not because she accepts it.

\paragraph{Idealized reflective endorsement.} ETH-01-a defines the good as what idealized reflective endorsement converges on: one outcome is better than another if every agent's idealized endorsement would prefer it. The good is an output of a procedure, the idealization of Definition~3.1 of ETH-01-a. ETH-01-a is explicit about the anti-realism about outcomes: where idealized agents would diverge, there is no fact about which outcome is better. That is the mark of pure procedural determination: the procedure's silence is not ignorance of an independent fact but the absence of one. The conditions of idealization, on the other hand, are not outputs of endorsement. ETH-01-a defends them against the objection that fully informed preferences lack authority by arguing for a weaker idealization, on which the agent's preference need only survive exposure to each relevant fact in turn. That argument concerns which procedure is correct, and it is the realist component. And the procedure is mandatory: the good, so defined, is what there is reason to prefer from the impartial standpoint, whatever an agent's actual endorsements.

The two cases also show how the present result fits MTE-02. The procedure of ETH-01-a is the same for every rational agent, and its outputs, the facts about what is better, may depend on the natures of those concerned. That is the structure of objectivity without absoluteness: a single mandatory procedure with nature-relative outputs.

\section{Objections}

\begin{objection}[Collapse into realism]
If constructivism is realist about its procedure, it is a form of realism. The realist holds that some normative facts are not constructed, and Corollary~\ref{c:global} concedes that.
\end{objection}

\begin{reply}
The realist the constructivist opposes holds that the outputs are true independently of the procedure, and that the procedure, if correct, is correct because it tracks them, as in perfect procedural justice. Theorem~\ref{t:main} denies this. The constructivist's procedure is correct because of its own features, such as the fairness of the conditions of choice, and its outputs are correct because they are its outputs, as in pure procedural justice. The two views differ in the order of explanation for every first-order moral truth, and agree only on there being one unconstructed normative fact. That is a difference large enough to make constructivism a distinct position, and small enough to make it stable.
\end{reply}

\begin{objection}[Constitutivism]
The procedure need not be mandatory in a realist sense, because it is constitutive of agency: anyone who acts at all is already committed to the categorical imperative, or to valuing her humanity \citep{Korsgaard1996}. The authority of the procedure is constructed from the nature of agency.
\end{objection}

\begin{reply}
That the procedure is constitutive of agency, if true, explains why agents are subject to it; it does not explain why they are required to accept its outputs. An agent may act defectively, and a creature may engage in something like agency without its constitutive commitments, as \citet{Enoch2006} argued with the case of the \enquote{shmagent}. The claim that one ought to be an agent rather than a shmagent, and so ought to accept the procedure's outputs, is a normative claim, and it is not an output of the procedure. Constitutivism is a theory of why the procedure is the correct one, and as such it is a candidate for the realist component, not a replacement for it.
\end{reply}

\begin{objection}[Euthyphro for procedures]
Is the procedure correct because it yields the right outcomes, or are the outcomes right because the procedure yields them? If the first, outcomes are prior and constructivism is false; if the second, any procedure would do.
\end{objection}

\begin{reply}
Neither. The outcomes are right because the procedure yields them, and the procedure is correct because of its own features, not because of its outcomes. This is the structure of every case of pure procedural justice: a lottery is fair because each ticket has an equal chance, not because it selects the right winner, and the winner is rightly selected because the lottery was fair. The second horn's claim that any procedure would do follows only if nothing but its outcomes could make a procedure correct, which is what the realism about procedures denies.
\end{reply}

\begin{objection}[Symmetry is too strong]
Assumption~\ref{a:symmetry} may fail: perhaps some procedures validate themselves in a way their rivals cannot, for instance because they are the procedures any validation must use.
\end{objection}

\begin{reply}
If some procedure is privileged in this way, its privilege is a fact about it that does not consist in its outputs, since its rivals have outputs of the same form; and that is realism about the procedure. The objection therefore leads to the conclusion of Theorem~\ref{t:main} by a different route. Assumption~\ref{a:symmetry} is used only in Lemma~\ref{l:external}, to show that the mandatoriness of the procedure cannot be fixed by chains of procedural endorsement in which every link is constructed.
\end{reply}

\section{What the result does not establish}

The theorem does not establish that constructivism is true of morality; it says what constructivism must be if it is. It does not identify the correct procedure, which is the task of the first-order theories that propose one; POL-01 and ETH-01 propose procedures for justice and for the good respectively, and the result is neutral between them and their rivals. It does not show that the realist component of constructivism is known by the same means as its outcomes; MTE-08 holds that moral truths are known by intuition, and the correctness of a procedure may be known that way. And it does not decide whether a single procedure underlies all of morality or different procedures underlie different parts of it, as the division of labour between POL-01 and ETH-01 suggests.

\section{Conclusion}

Constructivism constructs outcomes and lays out procedures. About the outcomes it is anti-realist: each moral truth holds because it is an output of the correct procedure, and where the procedure is silent there is no fact. About the procedure it is realist: that the procedure is the correct one, and that every rational agent is required to accept its outputs, is not the output of that procedure or of any other, and cannot be fixed by chains of procedural endorsement. The procedure itself is mandatory. Global constructivism is therefore incoherent, and a constructivism whose procedure varies with the agent's standpoint is relativism.

\begin{thebibliography}{99}\small
\bibitem[Enoch(2006)]{Enoch2006} Enoch, D. (2006). Agency, shmagency: Why normativity won't come from what is constitutive of action. \emph{The Philosophical Review}, 115(2), 169--198.
\bibitem[Enoch(2009)]{Enoch2009} Enoch, D. (2009). Can there be a global, interesting, coherent constructivism about practical reason? \emph{Philosophical Explorations}, 12(3), 319--339.
\bibitem[Fine(2012)]{Fine2012} Fine, K. (2012). Guide to ground. In F.~Correia and B.~Schnieder (Eds.), \emph{Metaphysical Grounding: Understanding the Structure of Reality} (pp.~37--80). Cambridge University Press.
\bibitem[Habermas(1990)]{Habermas1990} Habermas, J. (1990). \emph{Moral Consciousness and Communicative Action}. MIT Press.
\bibitem[Hussain and Shah(2006)]{HussainShah2006} Hussain, N. and Shah, N. (2006). Misunderstanding metaethics: Korsgaard's rejection of realism. In R.~Shafer-Landau (Ed.), \emph{Oxford Studies in Metaethics}, Vol.~1. Oxford University Press.
\bibitem[Korsgaard(1996)]{Korsgaard1996} Korsgaard, C.~M. (1996). \emph{The Sources of Normativity}. Cambridge University Press.
\bibitem[O'Neill(1989)]{ONeill1989} O'Neill, O. (1989). \emph{Constructions of Reason: Explorations of Kant's Practical Philosophy}. Cambridge University Press.
\bibitem[Rawls(1971)]{Rawls1971} Rawls, J. (1971). \emph{A Theory of Justice}. Harvard University Press.
\bibitem[Rawls(1980)]{Rawls1980} Rawls, J. (1980). Kantian constructivism in moral theory. \emph{The Journal of Philosophy}, 77(9), 515--572.
\bibitem[Rosen(2010)]{Rosen2010} Rosen, G. (2010). Metaphysical dependence: Grounding and reduction. In B.~Hale and A.~Hoffmann (Eds.), \emph{Modality: Metaphysics, Logic, and Epistemology} (pp.~109--135). Oxford University Press.
\bibitem[Scanlon(1998)]{Scanlon1998} Scanlon, T.~M. (1998). \emph{What We Owe to Each Other}. Harvard University Press.
\bibitem[Shafer-Landau(2003)]{ShaferLandau2003} Shafer-Landau, R. (2003). \emph{Moral Realism: A Defence}. Clarendon Press.
\bibitem[Street(2008)]{Street2008} Street, S. (2008). Constructivism about reasons. In R.~Shafer-Landau (Ed.), \emph{Oxford Studies in Metaethics}, Vol.~3. Oxford University Press.
\bibitem[Street(2010)]{Street2010} Street, S. (2010). What is constructivism in ethics and metaethics? \emph{Philosophy Compass}, 5(5), 363--384.
\end{thebibliography}

\end{document}
